
\documentclass{article}
\usepackage{colortbl}
\usepackage{makecell}
\usepackage{multirow}
\usepackage{supertabular}

\begin{document}

\newcounter{utterance}

\centering \large Interaction Transcript for game `codenames', experiment `frequency\_low', episode 0 with reflexion{-}Qwen3.5{-}9B.
\vspace{24pt}

{ \footnotesize  \setcounter{utterance}{1}
\setlength{\tabcolsep}{0pt}
\begin{supertabular}{c@{$\;$}|p{.15\linewidth}@{}p{.15\linewidth}p{.15\linewidth}p{.15\linewidth}p{.15\linewidth}p{.15\linewidth}}
    \# & $\;$A & \multicolumn{4}{c}{Game Master} & $\;\:$B\\
    \hline

    \theutterance \stepcounter{utterance}  
    & & & \multicolumn{2}{p{0.3\linewidth}}{
        \cellcolor[rgb]{0.9,0.9,0.9}{
            \makecell[{{p{\linewidth}}}]{
                \texttt{\tiny{[GM$|$GM]}}
                {'hidden': {'team': ['darer', 'cancelation'], 'innocent': ['ungodliness', 'astrobiology', 'gyration', 'fossilization'], 'opponent': [], 'assassin': ['caretaking']}, 'revealed': {'team': {'team': [], 'innocent': [], 'opponent': [], 'assassin': []}, 'opponent': {'team': [], 'innocent': [], 'opponent': [], 'assassin': []}}}
            }
        }
    }
    & & \\ \\

    \theutterance \stepcounter{utterance}  
    & & \multicolumn{4}{p{0.6\linewidth}}{
        \cellcolor[rgb]{0.9,0.9,0.9}{
            \makecell[{{p{\linewidth}}}]{
                \texttt{\tiny{[P1$\langle$GM]}}
                \texttt{You have attempted this same task before and did not succeed. These are the plans you wrote after those attempts.} \\
\texttt{Trial \#1: New plan: The previous attempt was aborted because the response included extra text or formatting errors outside the strict "CLUE: <WORD>" and "TARGETS: <TARGETS>" structure. I failed to adhere to the "never include any other text" rule. Next time, I will generate a response containing *only* the two required lines with no introductory phrases, explanations, or markdown formatting, ensuring the clue word is semantically related to the targets and does not contain parts of the target words themselves. For example, I will select "burn" for "candleholder" and "balsawood" and output strictly:} \\
\texttt{CLUE: burn} \\
\texttt{TARGETS: candleholder, balsawood} \\
\\ 
\texttt{Trial \#2: New plan: The previous attempt failed because the clue "dare" was a substring of the target word "darer," violating the rule that the clue cannot contain parts of the target words. Additionally, the clue "dare" is too specific and likely triggered the opponent's "undernourishment" (via "dare" {-}> "under" association? No, likely "dare" triggered "dangerousness" or similar in previous turns, but here it failed to guess "darer" directly while exposing "cancelation" or triggering a distractor). The core mistake was choosing a clue that is a prefix of a target word. Next time, I will select a clue that is semantically related but shares no letters or substrings with the targets. For "darer" and "cancelation," I will use "courage" (related to daring) and "void" (related to canceling), ensuring neither word contains "darer" or "cancelation" as substrings. I will strictly output only the two required lines.} \\
\texttt{CLUE: courage} \\
\texttt{TARGETS: darer, cancelation} \\
\\ 
\\ 
\texttt{Provide one single{-}word clue that relates to some of your team's words from the list below. You can choose to how many and to which words your clue relates to.} \\
\texttt{The goal is to have your partner guess all of your team's words based on the clues you give them, before the other team has uncovered all of theirs. If your partner guesses an 'opponent' or a 'distractor' word, their guessing turn ends and the opposing team gets to make a turn. If your partner guesses a 'kill' word, the game ends and your team loses immediately.} \\
\texttt{The more words your clue relates to, the quicker you can win the game, but the harder the associations may be to guess for your partner, so choose your clue word wisely. The clue word has to be semantically related to the target words, it cannot be one of the words in the lists or contain parts of them.} \\
\\ 
\texttt{Always give your single{-}word clue and your comma{-}separated list of related target words in the following format and make your answers as short as possible, never include any other text than is required in this form:} \\
\\ 
\texttt{CLUE: <WORD>} \\
\texttt{TARGETS: <TARGETS>} \\
\\ 
\texttt{Your team words are: darer, cancelation.} \\
\texttt{Your opponent's team words are: .} \\
\texttt{Distractor words are: ungodliness, astrobiology, gyration, fossilization.} \\
\texttt{The kill words are: caretaking.} \\
            }
        }
    }
    & \\ \\

    \theutterance \stepcounter{utterance}  
    & \multicolumn{4}{p{0.6\linewidth}}{
        \cellcolor[rgb]{0.9,0.9,0.9}{
            \makecell[{{p{\linewidth}}}]{
                \texttt{\tiny{[P1$\rangle$GM]}}
                \texttt{CLUE: courage} \\
\texttt{TARGETS: darer, cancelation} \\
            }
        }
    }
    & & \\ \\

    \theutterance \stepcounter{utterance}  
    & & & \multicolumn{2}{p{0.3\linewidth}}{
        \cellcolor[rgb]{0.9,0.9,0.9}{
            \makecell[{{p{\linewidth}}}]{
                \texttt{\tiny{[GM$|$GM]}}
                \texttt{courage} \\
            }
        }
    }
    & & \\ \\

    \theutterance \stepcounter{utterance}  
    & & & \multicolumn{2}{p{0.3\linewidth}}{
        \cellcolor[rgb]{0.9,0.9,0.9}{
            \makecell[{{p{\linewidth}}}]{
                \texttt{\tiny{[GM$|$GM]}}
                ['darer', 'cancelation']
            }
        }
    }
    & & \\ \\

    \theutterance \stepcounter{utterance}  
    & & & \multicolumn{2}{p{0.3\linewidth}}{
        \cellcolor[rgb]{0.9,0.9,0.9}{
            \makecell[{{p{\linewidth}}}]{
                \texttt{\tiny{[GM$|$GM]}}
                {'word': 'darer', 'assignment': 'team'}
            }
        }
    }
    & & \\ \\

    \theutterance \stepcounter{utterance}  
    & & & \multicolumn{2}{p{0.3\linewidth}}{
        \cellcolor[rgb]{0.9,0.9,0.9}{
            \makecell[{{p{\linewidth}}}]{
                \texttt{\tiny{[GM$|$GM]}}
                {'word': 'cancelation', 'assignment': 'team'}
            }
        }
    }
    & & \\ \\

    \theutterance \stepcounter{utterance}  
    & & & \multicolumn{2}{p{0.3\linewidth}}{
        \cellcolor[rgb]{0.9,0.9,0.9}{
            \makecell[{{p{\linewidth}}}]{
                \texttt{\tiny{[GM$|$GM]}}
                \texttt{opponent has won} \\
            }
        }
    }
    & & \\ \\

    \theutterance \stepcounter{utterance}  
    & \multicolumn{4}{p{0.6\linewidth}}{
        \cellcolor[rgb]{0.95,0.95,0.95}{
            \makecell[{{p{\linewidth}}}]{
                \texttt{\tiny{[P2$\rightarrow$P1]}}
                \texttt{Post{-}episode Reflexion (not sent to GM):} \\
\texttt{New plan: The previous attempt failed because the clue "courage" is semantically too broad and likely triggered the opponent's "distractor" words (e.g., "ungodliness" or "astrobiology") or caused the partner to guess a "kill" word like "caretaking" instead of the targets. A single clue must be specific enough to avoid ambiguity with distractors while remaining distinct from the targets. Next time, I will choose a clue that shares a very specific, unique semantic niche with both "darer" and "cancelation" but has no overlap with the distractor list. For instance, using "revoke" (specific to cancelation) might be too narrow, so I will try a concept like "undo" or "void" but verify it doesn't match distractors. However, since "courage" failed, I will pivot to a more abstract but precise concept like "risk" (linking to darer) paired with "nullify" (linking to cancelation) if allowed, or better yet, select a single word that bridges them uniquely, such as "rescind" (if it fits both) or choose a completely different angle. Actually, looking at the distractors, "courage" might have been too generic. I will select "daring" (too close to darer?) No. Let's try "void" for cancelation and see if it links to darer? No. The best strategy is to pick a word that is the *antonym* of the distractors or a very specific synonym. Let's try "retract" (cancelation) and see if it links to darer? No.} \\
\texttt{Wait, the goal is one single{-}word clue relating to *some* of the team's words. I can choose to relate to just one or both. If I relate to both, it's harder. Maybe relating to just "darer" with a specific word like "audacity" is safer? But the prompt says "relates to some of your team's words". If I pick "audacity", it relates to "darer". Does it relate to "cancelation"? No. That's fine. But did I fail because I related to both? The loss might be because "courage" triggered a kill word.} \\
\texttt{Let's re{-}evaluate: The clue "courage" is a synonym for "daring", which relates to "darer". It is also related to "cancelation"? Not really. Maybe the partner guessed "caretaking" because "courage" sounded like "care"? Or maybe "courage" triggered "ungodliness"?} \\
\texttt{Actually, the most likely error is that "courage" is too common and led to a "kill" word guess. I should pick a more obscure but precise word.} \\
\texttt{Let's try: CLUE: revoke (relates to cancelation). Does it relate to} \\
            }
        }
    }
    & & \\ \\

\end{supertabular}
}

\end{document}
